\documentclass[10pt]{article}
\usepackage{amsmath, amssymb, amsthm}
\usepackage{geometry}
\geometry{a4paper, margin=1in}

\usepackage{lmodern}
\usepackage{parskip}
\usepackage{enumitem}
\usepackage{hyperref}
\usepackage{xcolor}
\usepackage{fancyhdr}
\usepackage{graphicx}

\definecolor{mypink}{RGB}{255, 105, 180}
\definecolor{mydarkpink}{RGB}{199, 21, 133}
\definecolor{mygray}{RGB}{80, 80, 80}
\hypersetup{
    colorlinks=true,
    linkcolor=mydarkpink,
    urlcolor=mydarkpink,
    citecolor=mydarkpink
}

\pagestyle{fancy}
\fancyhf{}
\fancyhead[L]{\textcolor{mygray}{Elijah's Math Notes}}
\fancyhead[R]{\textcolor{mygray}{\thepage}}
\fancyfoot[C]{\textcolor{mygray}{Prepared by Elijah}}

\usepackage{titling}
\pretitle{\begin{center}\Large\color{mypink}}
\posttitle{\par\end{center}\vskip 0.5em}
\preauthor{\begin{center}\large\color{mygray}}
\postauthor{\par\end{center}}
\predate{\begin{center}\small\color{mygray}}
\postdate{\par\end{center}}

\theoremstyle{plain}
\newtheorem{theorem}{Theorem}[section]
\newtheorem{lemma}[theorem]{Lemma}
\newtheorem{proposition}[theorem]{Proposition}
\newtheorem{corollary}[theorem]{Corollary}

\theoremstyle{definition}
\newtheorem{definition}[theorem]{Definition}
\newtheorem{example}[theorem]{Example}

\theoremstyle{remark}
\newtheorem{remark}[theorem]{Remark}
\newtheorem{note}[theorem]{Note}

\newcommand{\R}{\mathbb{R}}
\newcommand{\N}{\mathbb{N}}
\newcommand{\Z}{\mathbb{Z}}
\newcommand{\todo}[1]{\textcolor{mydarkpink}{\textbf{TODO:} #1}}
\newcommand{\highlight}[1]{\textcolor{mypink}{#1}}

\title{Number Theory}
\author{Elijah Renner}
\date{\today}

\begin{document}

\maketitle
\tableofcontents

\section{Introduction}

Number theory studies the set of natural numbers \(\mathbb{N}=\{1,2,3,4,5,6,7,...\}\).

\(1+2+3+...+(n-1)+n=\frac{n(n+1)}{2}\).

Square numbers \(1,4,9,16,25,36,...\) are numbers that can be organized in squares.

Triangular numbers \(1,3,6,10,15,21,...\) are numbers that can be arranged in a triangle.

Prime numbers \(2,3,5,7,11,13,17,19,23,29,31,...\) are numbers whose divisors are only themselves and \(1\). Composite numbers are \(\mathbb{N}\) with the primes removed.

Perfect numbers \(6,28,496,...\) are numbers whose divisors sum to the perfect number.

Fibonacci numbers \(1,1,2,3,5,8,13,21,34,...\) are created by starting with 1 and 1. Then, each number in the list the sum of the previous two.

\section{Big O Notation}

Formally, we say that a function \( f(n) \) is \( O(g(n)) \) if there exist positive constants \( c \) and \( n_0 \) such that:
\[
f(n) \leq c \cdot g(n) \quad \text{for all } n \geq n_0.
\]

Here:
- \( f(n) \) represents the actual runtime or space complexity of the algorithm.
- \( g(n) \) is a comparison function that represents the upper bound.
- \( c \) and \( n_0 \) are constants.

\section*{Common Big O Notation Examples}

Here are some commonly used Big O notations:

\begin{itemize}
    \item \( O(1) \) Constant Time: The runtime does not depend on the size of the input.
    \item \( O(\log n) \) Logarithmic Time: The runtime grows logarithmically with the input size.
    \item \( O(n) \) Linear Time: The runtime grows linearly with the input size.
    \item \( O(n \log n) \) Quasi-Linear Time: The runtime grows at a rate proportional to \( n \log n \).
    \item \( O(n^2) \) Quadratic Time: The runtime grows quadratically with the input size.
    \item \( O(2^n) \) Exponential Time: The runtime doubles with each additional unit increase in input size.
    \item \( O(n!) \) Factorial Time: The runtime grows factorially with the input size.
\end{itemize}

\section*{Example Analysis}

Consider the function \( f(n) = 3n^2 + 5n + 2 \). To determine its Big O notation:

1. Identify the dominant term: In this case, \( 3n^2 \).
2. Drop the constant coefficients: The Big O notation becomes \( O(n^2) \).

Thus, \( f(n) = O(n^2) \).

\section*{Properties of Big O Notation}

Big O notation has some useful mathematical properties:
\begin{itemize}
    \item \textbf{Transitivity:} If \( f(n) = O(g(n)) \) and \( g(n) = O(h(n)) \), then \( f(n) = O(h(n)) \).
    \item \textbf{Addition:} If \( f_1(n) = O(g_1(n)) \) and \( f_2(n) = O(g_2(n)) \), then \( f_1(n) + f_2(n) = O(\max(g_1(n), g_2(n))) \).
    \item \textbf{Multiplication by a Constant:} If \( f(n) = O(g(n)) \), then \( c \cdot f(n) = O(g(n)) \) for any constant \( c > 0 \).
\end{itemize}

\end{document}